\section{Affine symmetries}\label{ar:sec}
\begingroup\sloppy\emergencystretch=\linewidth

This section extends the reduction of Theorem~\ref{sr:thm:main} from Euclidean isometries to affine symmetries. Its
central statement is a transfer principle (Theorem~\ref{ar:thm:transfer}). Let $N\ge3$ and $n$ be natural numbers, and
suppose that every finite set in general position with at least $N$ points that is mapped into itself by a Euclidean
isometry other than the identity contains $n$ points in convex position. Then the same holds with an injective affine
map $p\mapsto Lp+b$ in place of the isometry. The Euclidean statement is a \emph{hypothesis} of the transfer (a
parameter of the theorem), so the transfer is an implication and asserts neither statement on its own.
Corollary~\ref{ar:cor:transfer-forms} restates the transfer for Mathlib affine maps, for affine equivalences and for
maps given in coordinates. With $N=33$, $n=7$ and the Euclidean statement supplied by Theorem~\ref{sr:thm:main}, the
transfer gives five forms of the statement ``a set of at least $33$ points in general
position with a non-trivial affine symmetry contains $7$ points in convex position'' (Theorem~\ref{ar:thm:main}). Like
Theorem~\ref{sr:thm:main}, these five theorems take five of the six case classes of Definition~\ref{sr:def:classes}
(all except \lean{Fourfold32}) as hypotheses, so they are conditional (Remark~\ref{ar:rem:classes}).

The idea is as follows. The symmetry $g(p)=Lp+b$ permutes the finite set $S$, and $S$ contains three non-collinear
points, because it is in general position and has at least $N\ge3$ points. So some power $g^m$ with $m\ge1$ fixes these
three points, and then its linear part satisfies $L^m=1$. The averaged quadratic form $z\mapsto\sum_{k<m}\operatorname{nsq}(L^kz)$, where $\operatorname{nsq}$ is the squared
Euclidean norm, is $L$-invariant and positive definite. Writing it as $\operatorname{nsq}(Tz)$ with an explicit upper
triangular matrix $T$ (a Cholesky factor) turns $g$ into the Euclidean isometry $T\circ g\circ T^{-1}$ of the set
$T(S)$. Finally, $T$ preserves general position and cardinality, and $T^{-1}$ carries $n$ points of $T(S)$ in convex
position back to $n$ points of $S$ in convex position.

\paragraph{Sources.}
All declarations of this section are in the namespace \lean{ErdosSzekeres.SymReduction} of the file
\dirn{SYMLEAN/AffineReduction.lean} (SHA-256 \texttt{41cfc0b60aec}\dots), lines 1059--1533 (its \S\S9--10).
\begin{itemize}
\item Lines 1--1058 of that file are byte-identical to lines 1--1058 of \dirn{SYMLEAN/SymmetricReduction.lean}
(SHA-256 \texttt{6d5ec90a4215}\dots). Their declarations, among them \lean{eDist2}, \lean{IsEucIsometry},
\lean{Invariant}, the case classes and \lean{heptagon_of_nontrivial_isometry} (line 997 in both files), are treated in
the section of Theorem~\ref{sr:thm:main}; the theorems of this section use the copies in
\dirn{SYMLEAN/AffineReduction.lean}.
\item The repository definitions \lean{Point}, \lean{GeneralPosition}, \lean{ConvexPosition} and \lean{HasConvexNGon}
are in \dirn{ErdosSzekeres/Definitions.lean} (SHA-256 \texttt{2425500eefd1}\dots), and \lean{orient} is in
\dirn{ErdosSzekeres/Orientation.lean:9} (SHA-256 \texttt{f67955ac7091}\dots).
\item Two reformulations of these definitions are used. The first is \lean{convexPosition_iff}, at
\dirn{ErdosSzekeres/Basic.lean:15} (file SHA-256 \texttt{13667aa2d3dc}\dots). The second is
\lean{generalPosition_iff_orient_ne_zero}, at \dirn{ErdosSzekeres/OrientationGeometry.lean:66} (file SHA-256
\texttt{c7db9bbbb65e}\dots).
\item Toolchain: Lean 4.24.0, Mathlib \texttt{f897ebcf}.
\end{itemize}

\paragraph{Rigor and axioms.}
Every statement below that carries a Lean reference is \emph{Lean-checked}. The evidence is as follows.
\begin{itemize}
\item The file \dirn{SYMLEAN/AffineReductionAxioms.lean} (SHA-256 \texttt{f6cb9372b174}\dots) consists of the 1533
lines of \dirn{SYMLEAN/AffineReduction.lean}, followed by a blank line, nine \texttt{\#print axioms} lines and one
\texttt{\#check} line.
\item The run table \dirn{SYMLEAN/logs/runs.tsv} records that \dirn{SYMLEAN/AffineReduction.lean} and
\dirn{SYMLEAN/AffineReductionAxioms.lean} were each compiled on 1 November 2026 with exit code 0.
\item The recorded output \dirn{SYMLEAN/logs/AffineReductionAxioms.out} (SHA-256 \texttt{38d0b6f59398}\dots) of the
axioms file states
that each of the following nine theorems depends on exactly the axioms \texttt{propext}, \texttt{Classical.choice} and
\texttt{Quot.sound}: \lean{heptagon_of_nontrivial_affineEquiv}, \lean{heptagon_of_nontrivial_affineEquiv_image},
\lean{heptagon_of_nontrivial_affineMap}, \lean{heptagon_of_nontrivial_affine},
\lean{heptagon_of_nontrivial_affine_coords}, \lean{hasConvexNGon_of_affine_symmetry},
\lean{generalPosition_image_affineEquiv}, \lean{hasConvexNGon_image_affine} and
\lean{heptagon_of_nontrivial_isometry}.
\item Every other declaration of this section is used, directly or through other declarations, in the proof of one of
these nine theorems. So it depends on no further axiom. For \lean{lin2_fst} and \lean{lin2_snd}
(Lemma~\ref{ar:lem:lin2}(a)) this inference needs a supplement. Their Lean proofs are \lean{rfl}, and they are used only inside
\texttt{simp only} (\dirn{SYMLEAN/AffineReduction.lean}, lines 1180, 1187, 1193 and 1197), which may apply such lemmas as definitional rewrites without
recording them in the proof term. So they need not occur in the proof terms from which \texttt{\#print axioms}
collects the axioms of the nine theorems. Apart from equality and \lean{rfl}, which use no axiom, their statements and
proofs mention only \lean{lin2} applied to points and the coordinates, additions and multiplications in the definition
of \lean{lin2}. The axioms check of \lean{hasConvexNGon_of_affine_symmetry} reaches all of these, because its proof applies
the map $T$ of Proposition~\ref{ar:prop:averaging}(b), a value of \lean{lin2}, to points (Lean lines 1373--1376), and the
check also follows definitions. So these two lemmas depend on no
further axiom either.
\item The compilation log \dirn{SYMLEAN/logs/AffineReduction.out} of \dirn{SYMLEAN/AffineReduction.lean} is empty
(no message). So no declaration of the file uses \lean{sorry}.
\item The five case classes are definitions of type \lean{Prop} that occur as explicit premises, not axioms. The
output of the \texttt{\#check} line, in the same log as the axioms, displays them as premises of
\lean{heptagon_of_nontrivial_affineEquiv}.
\item No proof uses \lean{decide} or \lean{native_decide}. Nothing in this section is a computation.
\end{itemize}
The proofs below are \emph{written proofs (not refereed)}. They follow the Lean proofs step by step unless a proof says
otherwise. Statements tagged ``written step'' are not Lean declarations.

\paragraph{Dependencies.}
From the foundations: Definitions~\ref{fd:def:point}, \ref{fd:def:orient}, \ref{fd:def:gp}, \ref{fd:def:cp},
\ref{fd:def:ngon} and~\ref{fd:def:inv}; Lemma~\ref{fd:lem:collinear} (general position in terms of $\orient$); and
Lemma~\ref{fd:lem:cp-hull} (convex position in terms of convex hulls). The lemmas below re-prove, in the exact form
used here, the facts summarised in Lemmas~\ref{fd:lem:affine-orient} and~\ref{fd:lem:affine-cp}. From the reduction
for isometries: Definition~\ref{sr:def:iso} (\lean{eDist2} and \lean{IsEucIsometry}),
Definition~\ref{sr:def:classes} (the case classes) and Theorem~\ref{sr:thm:main}.

\subsection{Conventions}

\paragraph{Points, matrices, iterates.}
Points are elements of $\R^2$ (Lean \lean{Point}, the functions from the index type \lean{Fin}~$2$ to the reals),
written
$p=(p_0,p_1)$. Points and vectors are the same objects, and $p\pm q$ and $\lambda p$ are taken coordinatewise.
\begin{itemize}
\item A \emph{matrix} is a real $2\times2$ matrix $M=(M_{ij})_{i,j\in\{0,1\}}$ (Lean: a
\lean{Matrix} with rows and columns indexed by \lean{Fin}~$2$ and real entries).
\item $Mz$ is the matrix--vector product, $(Mz)_i=\sum_{j\in\{0,1\}}M_{ij}z_j$ (Lean \lean{Matrix.mulVec}). It is
linear in $z$, and $M(M'z)=(MM')z$ for matrices $M,M'$.
\item $M^k$ is the $k$-th matrix power, with $M^0=1$, the identity matrix, and $M^{k+1}=M^kM=MM^k$.
\item $\det M=M_{00}M_{11}-M_{01}M_{10}$ (Mathlib \lean{Matrix.det_fin_two}).
\item For a map $g:\R^2\to\R^2$, $g^k$ is the $k$-fold iterate (Lean \texttt{g\textasciicircum[k]}, from \lean{Nat.iterate}): $g^0=\mathrm{id}$ and
$g^{k+1}=g\circ g^k$. (The definition of \lean{Nat.iterate} recurses on the other side, $g^{k+1}=g^k\circ g$. The
form used here, $g^{k+1}(p)=g(g^k(p))$, is the lemma \lean{Function.iterate_succ_apply'}, which the Lean proof of
Lemma~\ref{ar:lem:iter} uses.)
\item ``$g\ne\mathrm{id}$'' is always an inequality of maps of the whole plane.
\end{itemize}

\paragraph{Finite sets.}
A finite set is a Lean \lean{Finset} of points, and $\card S$ is its cardinality. For a finite set $S$ and a map $f$,
$f(S)=\{f(p):p\in S\}$ is again a finite set (Lean \lean{Finset.image}); if $f$ is injective, then $\card{f(S)}=\card S$.
``$f$ maps $S$ into $S$'' means \lean{Invariant}~$f$~$S$: $f(p)\in S$ for all $p\in S$ (Definition~\ref{fd:def:inv};
Lean \dirn{SYMLEAN/AffineReduction.lean:271}). General position, convex position and \lean{HasConvexNGon} are as in
Definitions~\ref{fd:def:gp}, \ref{fd:def:cp} and~\ref{fd:def:ngon}. By Lemma~\ref{fd:lem:collinear}, a finite set $S$
is in general position if and only if $\orient(a,b,c)\ne0$ for all pairwise distinct $a,b,c\in S$. By
Lemma~\ref{fd:lem:cp-hull}, $S$ is in convex position if and only if $p\notin\conv(S\setminus\{p\})$ for every $p\in
S$. Here $\conv(A)$ is the convex hull, the smallest convex set containing $A$. Finally, \lean{HasConvexNGon}~$n$~$S$
means that some $U\subseteq S$ with $\card U=n$ is in convex position.

\paragraph{Euclidean isometries.}
$\operatorname{eDist2}(p,q)=(p_0-q_0)^2+(p_1-q_1)^2$, and a map $h:\R^2\to\R^2$ is a \emph{Euclidean isometry}
(\lean{IsEucIsometry}~$h$) if $\operatorname{eDist2}(h(p),h(q))=\operatorname{eDist2}(p,q)$ for \emph{all}
$p,q\in\R^2$ (Definition~\ref{sr:def:iso}; Lean \dirn{SYMLEAN/AffineReduction.lean:265} and~\dirn{:268}). The squared
distance is written out because the metric that Mathlib puts on \lean{Point} is the sup metric, not the Euclidean
one.

\paragraph{Affine maps (Mathlib).}
\begin{itemize}
\item An \emph{affine map} $f$ of $\R^2$ (Mathlib \lean{AffineMap}) is a map $f:\R^2\to\R^2$ together with a linear
map $f_{\mathrm{lin}}:\R^2\to\R^2$ (Lean \lean{f.linear}) such that $f(v+p)=f_{\mathrm{lin}}(v)+f(p)$ for all
$v,p\in\R^2$. The same letter denotes the map. A linear map $\varphi$ is an affine map with
$\varphi_{\mathrm{lin}}=\varphi$ (Lean \lean{LinearMap.toAffineMap}).
\item An \emph{affine equivalence} $f$ of $\R^2$ (Mathlib \lean{AffineEquiv}) is a bijection $f$ of $\R^2$ together
with a linear bijection $f_{\mathrm{lin}}$ with the same property. Its underlying affine map
(\lean{AffineEquiv.toAffineMap}) has the same map and the same linear part.
\item $\mathrm{refl}$ denotes the identity affine equivalence of $\R^2$, whose map is $\mathrm{id}$ (Lean
\lean{AffineEquiv.refl}, for the reals and \lean{Point}).
\item Two affine equivalences with the same map are equal (\lean{AffineEquiv.ext}), because the linear part is
determined by the map: $f_{\mathrm{lin}}(v)=f(v)-f(0)$, by the defining property with $p=0$. Hence, for an affine
equivalence $g$, $g\ne\mathrm{refl}$ holds if and only if $g\ne\mathrm{id}$ as maps.
\item For a linear map $\varphi$ of $\R^2$, $M_\varphi$ is the matrix whose $j$-th column is $\varphi(e_j)$, where
$e_0=(1,0)$ and $e_1=(0,1)$ (Lean \lean{LinearMap.toMatrix'}). Then $M_\varphi z=\varphi(z)$ for every $z$, since
$\varphi(z)=z_0\varphi(e_0)+z_1\varphi(e_1)$ (\lean{LinearMap.toMatrix'_mulVec}). The matrix of a composition is the
product of the matrices.
\end{itemize}

\subsection{The averaged quadratic form}

\begin{definition}[squared norm; \leanref{ErdosSzekeres.SymReduction.nsq}{SYMLEAN/AffineReduction.lean:1088}]\label{ar:def:nsq}
For $v\in\R^2$, $\operatorname{nsq}(v)=v_0^2+v_1^2$.
\end{definition}

\begin{lemma}[coordinates of $Az$; \leanref{ErdosSzekeres.SymReduction.mulVec2_fst}{SYMLEAN/AffineReduction.lean:1097}, \leanref{ErdosSzekeres.SymReduction.mulVec2_snd}{SYMLEAN/AffineReduction.lean:1101}]\label{ar:lem:mulvec}
For every matrix $A$ and every $z\in\R^2$: $(Az)_0=A_{00}z_0+A_{01}z_1$ \textup{(}\lean{mulVec2_fst}\textup{)} and
$(Az)_1=A_{10}z_0+A_{11}z_1$ \textup{(}\lean{mulVec2_snd}\textup{)}.
\end{lemma}

\begin{proof}
This is the definition $(Az)_i=\sum_{j\in\{0,1\}}A_{ij}z_j$, written out for $i=0$ and $i=1$.
\end{proof}

\begin{lemma}[elementary facts on $\operatorname{nsq}$; \leanref{ErdosSzekeres.SymReduction.nsq_nonneg}{SYMLEAN/AffineReduction.lean:1090}, \leanref{ErdosSzekeres.SymReduction.eDist2_eq_nsq}{SYMLEAN/AffineReduction.lean:1094}]\label{ar:lem:nsq}
The following hold.
\begin{enumerate}[label=\textup{(\alph*)}]
\item \textup{(}\lean{nsq_nonneg}\textup{)} $0\le\operatorname{nsq}(v)$ for every $v\in\R^2$.
\item \textup{(}\lean{eDist2_eq_nsq}\textup{)} $\operatorname{eDist2}(p,q)=\operatorname{nsq}(p-q)$ for all
$p,q\in\R^2$.
\end{enumerate}
\end{lemma}

\begin{proof}
(a) $\operatorname{nsq}(v)$ is a sum of two squares. (b) $(p-q)_i=p_i-q_i$, so
$\operatorname{nsq}(p-q)=(p_0-q_0)^2+(p_1-q_1)^2=\operatorname{eDist2}(p,q)$.
\end{proof}

\begin{definition}[Gram entries; \leanref{ErdosSzekeres.SymReduction.gA}{SYMLEAN/AffineReduction.lean:1106}, \leanref{ErdosSzekeres.SymReduction.gB}{SYMLEAN/AffineReduction.lean:1107}, \leanref{ErdosSzekeres.SymReduction.gC}{SYMLEAN/AffineReduction.lean:1108}]\label{ar:def:gram}
For a matrix $A$,
\[
g_A(A)=A_{00}^2+A_{10}^2,\qquad g_B(A)=A_{00}A_{01}+A_{10}A_{11},\qquad g_C(A)=A_{01}^2+A_{11}^2
\]
(Lean \lean{gA}, \lean{gB}, \lean{gC}). These are the entries of the Gram matrix of the two columns of $A$.
\end{definition}

\begin{lemma}[squared norm of $Az$; \leanref{ErdosSzekeres.SymReduction.nsq_mulVec}{SYMLEAN/AffineReduction.lean:1110}]\label{ar:lem:nsq-mulvec}
For every matrix $A$ and every $z\in\R^2$,
$\operatorname{nsq}(Az)=g_A(A)\,z_0^2+2\,g_B(A)\,z_0z_1+g_C(A)\,z_1^2$.
\end{lemma}

\begin{proof}
By Lemma~\ref{ar:lem:mulvec}, $\operatorname{nsq}(Az)=(A_{00}z_0+A_{01}z_1)^2+(A_{10}z_0+A_{11}z_1)^2$. Expanding gives
$(A_{00}^2+A_{10}^2)z_0^2+2(A_{00}A_{01}+A_{10}A_{11})z_0z_1+(A_{01}^2+A_{11}^2)z_1^2$.
\end{proof}

\begin{definition}[the averaged form; \leanref{ErdosSzekeres.SymReduction.avgQ}{SYMLEAN/AffineReduction.lean:1116}, \leanref{ErdosSzekeres.SymReduction.qA}{SYMLEAN/AffineReduction.lean:1120}, \leanref{ErdosSzekeres.SymReduction.qB}{SYMLEAN/AffineReduction.lean:1121}, \leanref{ErdosSzekeres.SymReduction.qC}{SYMLEAN/AffineReduction.lean:1122}]\label{ar:def:avgq}
For a matrix $L$, $n\in\N$ and $z\in\R^2$,
\[
\operatorname{avgQ}(L,n,z)=\sum_{k=0}^{n-1}\operatorname{nsq}(L^kz),
\]
\[
q_A(L,n)=\sum_{k=0}^{n-1}g_A(L^k),\qquad q_B(L,n)=\sum_{k=0}^{n-1}g_B(L^k),\qquad q_C(L,n)=\sum_{k=0}^{n-1}g_C(L^k)
\]
(Lean \lean{avgQ}, \lean{qA}, \lean{qB}, \lean{qC}; the sums run over \lean{Finset.range}~$n$, and they are $0$ for
$n=0$).
\end{definition}

\begin{lemma}[coefficients of the averaged form; \leanref{ErdosSzekeres.SymReduction.avgQ_eq}{SYMLEAN/AffineReduction.lean:1124}]\label{ar:lem:avgq-eq}
For every matrix $L$, every $n\in\N$ and every $z\in\R^2$,
\[
\operatorname{avgQ}(L,n,z)=q_A(L,n)\,z_0^2+2\,q_B(L,n)\,z_0z_1+q_C(L,n)\,z_1^2 .
\]
\end{lemma}

\begin{proof}
Apply Lemma~\ref{ar:lem:nsq-mulvec} to $A=L^k$ for each $k<n$ and add the $n$ equalities. On the left this gives
$\operatorname{avgQ}(L,n,z)$. On the right, the coefficients of $z_0^2$, $2z_0z_1$ and $z_1^2$ add up to $q_A(L,n)$,
$q_B(L,n)$ and $q_C(L,n)$. (Lean argues by induction on $n$; each induction step is the addition of one such
equality.)
\end{proof}

\begin{lemma}[positivity of the averaged form; \leanref{ErdosSzekeres.SymReduction.avgQ_ge}{SYMLEAN/AffineReduction.lean:1133}, \leanref{ErdosSzekeres.SymReduction.qA_pos}{SYMLEAN/AffineReduction.lean:1152}, \leanref{ErdosSzekeres.SymReduction.qdisc_pos}{SYMLEAN/AffineReduction.lean:1158}]\label{ar:lem:avgq-pos}
Let $L$ be a matrix and $n\in\N$ with $0<n$. Then:
\begin{enumerate}[label=\textup{(\alph*)}]
\item \textup{(}\lean{avgQ_ge}\textup{)} $\operatorname{nsq}(z)\le\operatorname{avgQ}(L,n,z)$ for every $z\in\R^2$;
\item \textup{(}\lean{qA_pos}\textup{)} $0<q_A(L,n)$;
\item \textup{(}\lean{qdisc_pos}\textup{)} $0<q_A(L,n)\,q_C(L,n)-q_B(L,n)^2$.
\end{enumerate}
\end{lemma}

\begin{proof}
(a) Since $n>0$, the sum defining $\operatorname{avgQ}(L,n,z)$ contains the term $k=0$, which is
$\operatorname{nsq}(L^0z)=\operatorname{nsq}(z)$. The other terms are $\ge0$ by Lemma~\ref{ar:lem:nsq}(a).

Write $\alpha=q_A(L,n)$, $\beta=q_B(L,n)$, $\gamma=q_C(L,n)$.

(b) Apply (a) to $z=(1,0)$ and use Lemma~\ref{ar:lem:avgq-eq}: $1=\operatorname{nsq}(1,0)\le\operatorname{avgQ}(L,n,(1,0))
=\alpha$. Hence $\alpha\ge1>0$.

(c) Apply (a) to $z=(-\beta,\alpha)$ and use Lemma~\ref{ar:lem:avgq-eq}:
\[
\alpha^2+\beta^2=\operatorname{nsq}(-\beta,\alpha)\le\alpha\beta^2+2\beta(-\beta)\alpha+\gamma\alpha^2
=\alpha(\alpha\gamma-\beta^2).
\]
Hence $\alpha(\alpha\gamma-\beta^2)\ge\alpha^2+\beta^2\ge\alpha^2>0$, where the last step uses (b). Since $\alpha>0$, we
get $\alpha\gamma-\beta^2>0$.
\end{proof}

\begin{lemma}[invariance of the averaged form; \leanref{ErdosSzekeres.SymReduction.avgQ_mulVec}{SYMLEAN/AffineReduction.lean:1139}]\label{ar:lem:avgq-inv}
Let $L$ be a matrix and $n\in\N$ with $L^n=1$. Then $\operatorname{avgQ}(L,n,Lz)=\operatorname{avgQ}(L,n,z)$ for every
$z\in\R^2$.
\end{lemma}

\begin{proof}
Put $f(k)=\operatorname{nsq}(L^kz)$ for $k\in\N$. Since $L^k(Lz)=(L^kL)z=L^{k+1}z$, we have
$\operatorname{avgQ}(L,n,Lz)=\sum_{k=0}^{n-1}f(k+1)$. Splitting off the last term, respectively the first term, of
$\sum_{k=0}^{n}f(k)$ gives
\[
\sum_{k=0}^{n-1}f(k)+f(n)=\sum_{k=0}^{n}f(k)=\sum_{k=0}^{n-1}f(k+1)+f(0).
\]
Here $f(n)=\operatorname{nsq}(L^nz)=\operatorname{nsq}(z)=\operatorname{nsq}(L^0z)=f(0)$, because $L^n=1=L^0$. Hence
$\sum_{k=0}^{n-1}f(k+1)=\sum_{k=0}^{n-1}f(k)=\operatorname{avgQ}(L,n,z)$. No positivity of $n$ is needed; for $n=0$
both sides are $0$.
\end{proof}

\subsection{Linear maps with explicit matrices; the Cholesky factor}

\begin{definition}[the linear map $\operatorname{lin2}$; \leanref{ErdosSzekeres.SymReduction.lin2}{SYMLEAN/AffineReduction.lean:1167}]\label{ar:def:lin2}
For $a,b,c,d\in\R$, $\operatorname{lin2}(a,b,c,d):\R^2\to\R^2$ is the $\R$-linear map
$p\mapsto(ap_0+bp_1,\ cp_0+dp_1)$, that is, multiplication by the matrix
$\begin{pmatrix}a&b\\c&d\end{pmatrix}$. In Lean it is a bundled linear map; its additivity and homogeneity are
discharged inside the definition by coordinate computations.
\end{definition}

\begin{lemma}[coordinates and orientation for $\operatorname{lin2}$; \leanref{ErdosSzekeres.SymReduction.lin2_fst}{SYMLEAN/AffineReduction.lean:1174}, \leanref{ErdosSzekeres.SymReduction.lin2_snd}{SYMLEAN/AffineReduction.lean:1175}, \leanref{ErdosSzekeres.SymReduction.orient_lin2}{SYMLEAN/AffineReduction.lean:1177}]\label{ar:lem:lin2}
Let $a,b,c,d\in\R$ and $\Lambda=\operatorname{lin2}(a,b,c,d)$.
\begin{enumerate}[label=\textup{(\alph*)}]
\item \textup{(}\lean{lin2_fst}, \lean{lin2_snd}\textup{)} $(\Lambda p)_0=ap_0+bp_1$ and $(\Lambda p)_1=cp_0+dp_1$ for
every $p\in\R^2$.
\item \textup{(}\lean{orient_lin2}\textup{)} $\orient(\Lambda x,\Lambda y,\Lambda z)=(ad-bc)\orient(x,y,z)$ for all
$x,y,z\in\R^2$.
\end{enumerate}
\end{lemma}

\begin{proof}
(a) This is the definition. (b) By definition, $\orient(X,Y,Z)=(Y-X)_0(Z-X)_1-(Y-X)_1(Z-X)_0$. Put $u=y-x$ and $w=z-x$.
By linearity, $\Lambda y-\Lambda x=\Lambda u$ and $\Lambda z-\Lambda x=\Lambda w$. Hence, by (a),
\[
\orient(\Lambda x,\Lambda y,\Lambda z)=(au_0+bu_1)(cw_0+dw_1)-(cu_0+du_1)(aw_0+bw_1)=(ad-bc)(u_0w_1-u_1w_0),
\]
and $u_0w_1-u_1w_0=\orient(x,y,z)$. This is the case $v=0$ of Lemma~\ref{fd:lem:affine-orient}.
\end{proof}

\begin{lemma}[the Cholesky factor and its inverse; \leanref{ErdosSzekeres.SymReduction.nsq_lin2_chol}{SYMLEAN/AffineReduction.lean:1184}, \leanref{ErdosSzekeres.SymReduction.lin2_chol_inv_left}{SYMLEAN/AffineReduction.lean:1191}, \leanref{ErdosSzekeres.SymReduction.lin2_chol_inv_right}{SYMLEAN/AffineReduction.lean:1195}]\label{ar:lem:chol}
Let $s,t,\beta\in\R$. Put $T=\operatorname{lin2}(s,\beta/s,0,t)$ and $T'=\operatorname{lin2}(1/s,\,-\beta/(s^2t),\,0,\,1/t)$.
\textup(Division follows the Lean convention $x/0=0$, so $T$ and $T'$ are defined for all $s,t,\beta$; each part
below states the non-vanishing it assumes.\textup)
\begin{enumerate}[label=\textup{(\alph*)}]
\item \textup{(}\lean{nsq_lin2_chol}\textup{)} If $s\ne0$, then for every $z\in\R^2$,
\[
\operatorname{nsq}(Tz)=s^2z_0^2+2\beta z_0z_1+\Bigl(\frac{\beta^2}{s^2}+t^2\Bigr)z_1^2 .
\]
\item \textup{(}\lean{lin2_chol_inv_left}\textup{)} If $s\ne0$ and $t\ne0$, then $T'(Tx)=x$ for every $x\in\R^2$.
\item \textup{(}\lean{lin2_chol_inv_right}\textup{)} If $s\ne0$ and $t\ne0$, then $T(T'x)=x$ for every $x\in\R^2$.
\end{enumerate}
\end{lemma}

\begin{proof}
By Lemma~\ref{ar:lem:lin2}(a), $Tz=(sz_0+(\beta/s)z_1,\ tz_1)$ and $T'x=(x_0/s-\beta x_1/(s^2t),\ x_1/t)$.

(a) $\operatorname{nsq}(Tz)=s^2z_0^2+2s(\beta/s)z_0z_1+(\beta/s)^2z_1^2+t^2z_1^2$, and $s(\beta/s)=\beta$ because
$s\ne0$.

(b) The first coordinate of $T'(Tx)$ is $\frac1s\bigl(sx_0+\frac\beta sx_1\bigr)-\frac{\beta}{s^2t}\,tx_1
=x_0+\frac{\beta x_1}{s^2}-\frac{\beta x_1}{s^2}=x_0$. The second is $\frac1t\,tx_1=x_1$.

(c) The first coordinate of $T(T'x)$ is $s\bigl(\frac{x_0}s-\frac{\beta x_1}{s^2t}\bigr)+\frac\beta s\cdot\frac{x_1}t
=x_0-\frac{\beta x_1}{st}+\frac{\beta x_1}{st}=x_0$. The second is $t\cdot\frac{x_1}t=x_1$.
\end{proof}

\subsection{Affine maps in matrix form}

In this subsection, $g:\R^2\to\R^2$ is a map, $L$ a matrix and $b\in\R^2$; ``$g(p)=Lp+b$ for all $p$'' is the Lean
hypothesis \lean{hg}.

\begin{lemma}[differences of iterates; \leanref{ErdosSzekeres.SymReduction.iterate_sub_of_affine}{SYMLEAN/AffineReduction.lean:1199}]\label{ar:lem:iter}
If $g(p)=Lp+b$ for all $p\in\R^2$, then $g^k(p)-g^k(q)=L^k(p-q)$ for all $k\in\N$ and $p,q\in\R^2$.
\end{lemma}

\begin{proof}
Induction on $k$. For $k=0$ both sides equal $p-q$, as $L^0=1$. If the claim holds for $k$, then
\[
g^{k+1}(p)-g^{k+1}(q)=\bigl(Lg^k(p)+b\bigr)-\bigl(Lg^k(q)+b\bigr)=L\bigl(g^k(p)-g^k(q)\bigr)=L\bigl(L^k(p-q)\bigr)
=L^{k+1}(p-q).\qedhere
\]
\end{proof}

\begin{lemma}[periodic points; \leanref{ErdosSzekeres.SymReduction.exists_periodic_of_injective}{SYMLEAN/AffineReduction.lean:1208}]\label{ar:lem:periodic}
Let $g:\R^2\to\R^2$ be injective, let $S$ be a finite set that $g$ maps into itself, and let $q\in S$. Then there is
$m\in\N$ with $0<m$ and $g^m(q)=q$.
\end{lemma}

\begin{proof}
By induction on $i$, $g^i(q)\in S$ for every $i\in\N$: $g^0(q)=q\in S$, and $g^{i+1}(q)=g(g^i(q))\in S$ because $g$
maps $S$ into itself. (In Lean this is \lean{iterate_mem}, \dirn{SYMLEAN/AffineReduction.lean:69}.) Since $\N$ is
infinite and $S$ is finite, the map $i\mapsto g^i(q)$ from $\N$ to $S$ is not injective (pigeonhole). So there are
$i<j$ with $g^i(q)=g^j(q)$. Put $m=j-i>0$. Then $g^i(g^m(q))=g^{i+m}(q)=g^j(q)=g^i(q)$. The map $g^i$ is injective, as
a composition of injective maps, so $g^m(q)=q$.
\end{proof}

\begin{lemma}[a matrix fixing two independent vectors; \leanref{ErdosSzekeres.SymReduction.matrix_eq_one_of_fix}{SYMLEAN/AffineReduction.lean:1217}]\label{ar:lem:fix}
Let $M$ be a matrix and $u,w\in\R^2$ with $u_0w_1-u_1w_0\ne0$, $Mu=u$ and $Mw=w$. Then $M=1$.
\end{lemma}

\begin{proof}
Put $\delta=u_0w_1-u_1w_0$. By Lemma~\ref{ar:lem:mulvec}, the hypotheses say that the four quantities
\[
E_1=M_{00}u_0+M_{01}u_1-u_0,\quad E_2=M_{10}u_0+M_{11}u_1-u_1,\quad
E_3=M_{00}w_0+M_{01}w_1-w_0,\quad E_4=M_{10}w_0+M_{11}w_1-w_1
\]
vanish. Expanding gives the polynomial identities
\[
w_1E_1-u_1E_3=(M_{00}-1)\delta,\quad -w_0E_1+u_0E_3=M_{01}\delta,\quad
w_1E_2-u_1E_4=M_{10}\delta,\quad -w_0E_2+u_0E_4=(M_{11}-1)\delta .
\]
So all four right-hand sides vanish. As $\delta\ne0$, we get $M_{00}=1$, $M_{01}=0$, $M_{10}=0$ and $M_{11}=1$, that
is, $M=1$.
\end{proof}

\begin{lemma}[orientation under affine maps; \leanref{ErdosSzekeres.SymReduction.orient_affine}{SYMLEAN/AffineReduction.lean:1241}]\label{ar:lem:orient-affine}
If $g(p)=Lp+b$ for all $p\in\R^2$, then $\orient(g(x),g(y),g(z))=\det(L)\,\orient(x,y,z)$ for all $x,y,z\in\R^2$.
\end{lemma}

\begin{proof}
We have $g(y)-g(x)=L(y-x)$ and $g(z)-g(x)=L(z-x)$, since the translation $b$ cancels. So the computation in the proof
of Lemma~\ref{ar:lem:lin2}(b), with the entries $L_{00},L_{01},L_{10},L_{11}$ in place of the four parameters of
$\operatorname{lin2}$, gives the factor $L_{00}L_{11}-L_{01}L_{10}=\det L$. This is Lemma~\ref{fd:lem:affine-orient}.
\end{proof}

\begin{lemma}[injectivity; \leanref{ErdosSzekeres.SymReduction.injective_of_affine}{SYMLEAN/AffineReduction.lean:1305}]\label{ar:lem:inj}
If $g(p)=Lp+b$ for all $p\in\R^2$ and $\det L\ne0$, then $g$ is injective.
\end{lemma}

\begin{proof}
Let $g(p)=g(q)$ and put $v=p-q$. Then $Lv=Lp-Lq=(Lp+b)-(Lq+b)=g(p)-g(q)=0$. So
$h_0:=L_{00}v_0+L_{01}v_1=0$ and $h_1:=L_{10}v_0+L_{11}v_1=0$ by Lemma~\ref{ar:lem:mulvec}. Then
\[
(\det L)\,v_0=L_{11}h_0-L_{01}h_1=0,\qquad (\det L)\,v_1=L_{00}h_1-L_{10}h_0=0,
\]
so $v=0$ because $\det L\ne0$. Hence $p=q$.
\end{proof}

\begin{lemma}[matrix form of a Mathlib affine map; \leanref{ErdosSzekeres.SymReduction.affineMap_eq_mulVec}{SYMLEAN/AffineReduction.lean:1284}]\label{ar:lem:matrix-form}
Let $f$ be an affine map of $\R^2$ \textup(Mathlib \lean{AffineMap}\textup) and let $M_f$ be the matrix of its linear
part $f_{\mathrm{lin}}$ \textup(Lean \lean{LinearMap.toMatrix'} applied to \lean{f.linear}\textup). Then $f(p)=M_fp+f(0)$ for every
$p\in\R^2$.
\end{lemma}

\begin{proof}
By the defining property of an affine map with $v=p$ and base point $0$, $f(p)=f(p+0)=f_{\mathrm{lin}}(p)+f(0)$ (in
Mathlib, \lean{AffineMap.decomp}). By the conventions, $f_{\mathrm{lin}}(p)=M_fp$.
\end{proof}

\subsection{Transport of general position and convex position}

\begin{lemma}[general position under orientation-scaling maps; \leanref{ErdosSzekeres.SymReduction.generalPosition_image_of_orient}{SYMLEAN/AffineReduction.lean:1249}]\label{ar:lem:gp-orient}
Let $f:\R^2\to\R^2$ be a map and $D\in\R$ with $D\ne0$ and $\orient(f(x),f(y),f(z))=D\,\orient(x,y,z)$ for all
$x,y,z\in\R^2$. If a finite set $S$ is in general position, then so is $f(S)$.
\end{lemma}

\begin{proof}
By Lemma~\ref{fd:lem:collinear} it suffices to show $\orient(X,Y,Z)\ne0$ for pairwise distinct $X,Y,Z\in f(S)$. Write
$X=f(x)$, $Y=f(y)$, $Z=f(z)$ with $x,y,z\in S$. Then $x,y,z$ are pairwise distinct: for instance $x=y$ would give
$X=Y$. By Lemma~\ref{fd:lem:collinear} again, $\orient(x,y,z)\ne0$, and so
$\orient(X,Y,Z)=D\,\orient(x,y,z)\ne0$. Injectivity of $f$ is not assumed and not needed.
\end{proof}

\begin{lemma}[general position under affine maps; \leanref{ErdosSzekeres.SymReduction.generalPosition_image_affine}{SYMLEAN/AffineReduction.lean:1278}, \leanref{ErdosSzekeres.SymReduction.generalPosition_image_affineEquiv}{SYMLEAN/AffineReduction.lean:1290}]\label{ar:lem:gp-affine}
The following hold.
\begin{enumerate}[label=\textup{(\alph*)}]
\item \textup{(}\lean{generalPosition_image_affine}\textup{)} Let $g:\R^2\to\R^2$, $L$ a matrix and $b\in\R^2$ with
$g(p)=Lp+b$ for all $p$, and $\det L\ne0$. If a finite set $S$ is in general position, then so is $g(S)$.
\item \textup{(}\lean{generalPosition_image_affineEquiv}\textup{)} If a finite set $S$ is in general position and $f$
is an affine equivalence of $\R^2$ \textup(Mathlib \lean{AffineEquiv}\textup), then $f(S)$ is in general position.
\end{enumerate}
\end{lemma}

\begin{proof}
(a) By Lemma~\ref{ar:lem:orient-affine}, the map $g$ satisfies the hypothesis of Lemma~\ref{ar:lem:gp-orient} with
$D=\det L\ne0$.

(b) By Lemma~\ref{ar:lem:matrix-form}, applied to the underlying affine map, $f(p)=Mp+f(0)$ for all $p$, where $M$ is
the matrix of $f_{\mathrm{lin}}$. The linear bijection $f_{\mathrm{lin}}$ has a linear inverse $\psi$. The product $MM_\psi$
is the matrix of $f_{\mathrm{lin}}\circ\psi=\mathrm{id}$, so $MM_\psi=1$ and $\det M\cdot\det M_\psi=1$; hence $\det
M\ne0$. (Lean uses \lean{LinearMap.det_toMatrix'} and \lean{LinearEquiv.isUnit_det'}.) Now apply (a).
\end{proof}

\begin{lemma}[convex position under injective affine maps; \leanref{ErdosSzekeres.SymReduction.convexPosition_image_affine}{SYMLEAN/AffineReduction.lean:1263}]\label{ar:lem:cp-affine}
Let $S$ be a finite set in convex position and let $f$ be an affine map of $\R^2$ \textup(Mathlib
\lean{AffineMap}\textup) that is injective. Then $f(S)$ is in convex position.
\end{lemma}

\begin{proof}
\emph{Step 1: $f(\conv A)=\conv(f(A))$ for every $A\subseteq\R^2$} (in Mathlib, \lean{AffineMap.image_convexHull}). By
Lemma~\ref{ar:lem:matrix-form}, $f(p)=Mp+c$ with $M=M_f$ and $c=f(0)$. So for $x,y\in\R^2$ and $\lambda\in[0,1]$,
\[
f\bigl((1-\lambda)x+\lambda y\bigr)=(1-\lambda)Mx+\lambda My+\bigl((1-\lambda)+\lambda\bigr)c
=(1-\lambda)f(x)+\lambda f(y).
\]
Hence images and preimages of convex sets under $f$ are convex. Let $C$ be convex. If $X=f(x)$ and $Y=f(y)$ with
$x,y\in C$, then $(1-\lambda)x+\lambda y\in C$, so $(1-\lambda)X+\lambda Y=f((1-\lambda)x+\lambda y)\in f(C)$. Let
$C'$ be convex. If $f(x),f(y)\in C'$, then $f((1-\lambda)x+\lambda y)=(1-\lambda)f(x)+\lambda f(y)\in C'$, so
$(1-\lambda)x+\lambda y\in f^{-1}(C')$. Now $f(\conv A)$ is a convex set containing $f(A)$, so
it contains $\conv(f(A))$. And $f^{-1}(\conv(f(A)))$ is a convex set containing $A$, so it contains $\conv A$; that is,
$f(\conv A)\subseteq\conv(f(A))$.

\emph{Step 2.} By Lemma~\ref{fd:lem:cp-hull} we must show $P\notin\conv(f(S)\setminus\{P\})$ for every $P\in f(S)$.
Write $P=f(q)$ with $q\in S$. Since $f$ is injective, $f(S)\setminus\{f(q)\}=f(S\setminus\{q\})$. Indeed, if $r\in S$
and $f(r)\ne f(q)$ then $r\ne q$; conversely, if $r\in S$ and $r\ne q$ then $f(r)\ne f(q)$. Suppose
$f(q)\in\conv(f(S)\setminus\{f(q)\})$. By this identity and Step~1, this set is
$\conv(f(S\setminus\{q\}))=f(\conv(S\setminus\{q\}))$. So $f(q)=f(r)$ for some $r\in\conv(S\setminus\{q\})$.
Injectivity gives $q=r\in\conv(S\setminus\{q\})$, which contradicts the convex position of $S$
(Lemma~\ref{fd:lem:cp-hull}).
\end{proof}

\begin{lemma}[convex $n$-gons under injective affine maps; \leanref{ErdosSzekeres.SymReduction.hasConvexNGon_image_affine}{SYMLEAN/AffineReduction.lean:1298}]\label{ar:lem:ngon-affine}
Let $n\in\N$, let $S$ be a finite set with \lean{HasConvexNGon}~$n$~$S$, and let $f$ be an injective affine map of $\R^2$
\textup(Mathlib \lean{AffineMap}\textup). Then \lean{HasConvexNGon}~$n$ holds for $f(S)$.
\end{lemma}

\begin{proof}
Take $U\subseteq S$ with $\card U=n$ in convex position. Then $f(U)\subseteq f(S)$, $\card{f(U)}=\card U=n$ since $f$ is
injective, and $f(U)$ is in convex position by Lemma~\ref{ar:lem:cp-affine}.
\end{proof}

\subsection{The transfer theorem}

\begin{definition}[the Euclidean statement $\mathcal E(N,n)$; written step]\label{ar:def:E}
For $N,n\in\N$, $\mathcal E(N,n)$ is the following proposition. \emph{For every finite set $S'\subseteq\R^2$ and every
map $h:\R^2\to\R^2$: if $S'$ is in general position, $N\le\card{S'}$, $h$ is a Euclidean isometry \textup(that is,
$\operatorname{eDist2}(h(p),h(q))=\operatorname{eDist2}(p,q)$ for all $p,q\in\R^2$\textup), $h\ne\mathrm{id}$, and
$h(p)\in S'$ for all $p\in S'$, then \lean{HasConvexNGon}~$n$~$S'$.} This is the type of the hypothesis \lean{hE} of
Theorem~\ref{ar:thm:transfer} and Corollary~\ref{ar:cor:transfer-forms}, up to two changes of notation: the bound
variables, \lean{S} and \lean{g} in Lean, are called $S'$ and $h$ here, and \lean{IsEucIsometry} and \lean{Invariant} are
unfolded. The quantifiers, the five premises and their order, and the conclusion are those of Lean. It has no name in
Lean.
\end{definition}

The next proposition collects the steps of the proof of Theorem~\ref{ar:thm:transfer} that construct the Euclidean
isometry. It is stated separately because the remark on affine symmetries in \dirn{WRITEUP/P1/p1.tex} (Remark
\texttt{rem:affine}) uses this construction.

\begin{proposition}[averaging an affine symmetry; written step]\label{ar:prop:averaging}
Let $S$ be a finite set in general position with $\card S\ge3$. Let $L$ be a matrix, $b\in\R^2$ and
$g:\R^2\to\R^2$ with $g(p)=Lp+b$ for all $p$. Assume that $g$ is injective and maps $S$ into itself.
\begin{enumerate}[label=\textup{(\alph*)}]
\item There is an integer $m\ge1$ with $L^m=1$ and $g^m=\mathrm{id}$.
\item Fix any integer $m\ge1$ with $L^m=1$, and put $\alpha=q_A(L,m)$, $\beta=q_B(L,m)$, $\gamma=q_C(L,m)$. Then
$\alpha>0$ and $\alpha\gamma-\beta^2>0$. Put
\[
s=\sqrt\alpha,\qquad t=\sqrt{(\alpha\gamma-\beta^2)/\alpha},\qquad
T=\operatorname{lin2}(s,\beta/s,0,t),\qquad T'=\operatorname{lin2}(1/s,\,-\beta/(s^2t),\,0,\,1/t).
\]
Then $s>0$ and $t>0$. The maps $T$ and $T'$ are mutually inverse linear bijections of $\R^2$. Moreover
$\orient(Tx,Ty,Tz)=st\cdot\orient(x,y,z)$ for all $x,y,z$, with $st>0$, and
$\operatorname{nsq}(Tz)=\operatorname{avgQ}(L,m,z)$ for every $z$.
\item With $m$, $T$, $T'$ as in \textup{(b)}, the map $h=T\circ g\circ T'$ is a Euclidean isometry, and $h$ maps
$T(S)$ into $T(S)$.
\item $h^k=T\circ g^k\circ T'$ for every $k\in\N$. Hence $h^k=\mathrm{id}$ if and only if $g^k=\mathrm{id}$. In
particular, $h\ne\mathrm{id}$ if and only if $g\ne\mathrm{id}$, and $h$ and $g$ have the same order (the least $k\ge1$
with $g^k=\mathrm{id}$, which exists by \textup{(a)}).
\end{enumerate}
\end{proposition}

\begin{proof}
(a) Since $\card S\ge3$, $S$ contains pairwise distinct points $p_0,p_1,p_2$. By general position and
Lemma~\ref{fd:lem:collinear}, $\delta:=\orient(p_0,p_1,p_2)\ne0$. By Lemma~\ref{ar:lem:periodic} there are
$m_0,m_1,m_2\ge1$ with $g^{m_i}(p_i)=p_i$. Put $m=m_0m_1m_2\ge1$. For each $i$, $m=r_im_i$ for an integer $r_i\ge1$. By
induction on $r$, $g^{rm_i}(p_i)=p_i$: for $r=0$ this is $g^0(p_i)=p_i$, and
$g^{(r+1)m_i}(p_i)=g^{m_i}(g^{rm_i}(p_i))=g^{m_i}(p_i)=p_i$. So $g^m(p_i)=p_i$ for $i=0,1,2$.

Put $u=p_1-p_0$ and $w=p_2-p_0$. By Lemma~\ref{ar:lem:iter}, $L^mu=g^m(p_1)-g^m(p_0)=p_1-p_0=u$, and likewise
$L^mw=w$. By the definition of $\orient$, $u_0w_1-u_1w_0=\orient(p_0,p_1,p_2)=\delta\ne0$. So Lemma~\ref{ar:lem:fix}
gives $L^m=1$.

Finally, for every $p\in\R^2$, Lemma~\ref{ar:lem:iter} gives $g^m(p)-g^m(p_0)=L^m(p-p_0)=p-p_0$. Since
$g^m(p_0)=p_0$, this gives $g^m(p)=p$. So $g^m=\mathrm{id}$. (The Lean proof derives $L^m=1$ in the same way, as its
step \lean{hLm}. It does not derive or use $g^m=\mathrm{id}$.)

(b) By Lemma~\ref{ar:lem:avgq-pos}(b),(c) with $n=m\ge1$, $\alpha>0$ and $\alpha\gamma-\beta^2>0$. So
$(\alpha\gamma-\beta^2)/\alpha>0$; hence $s>0$, $t>0$, $s^2=\alpha$ and $t^2=(\alpha\gamma-\beta^2)/\alpha$. As $s\ne0$
and $t\ne0$, Lemma~\ref{ar:lem:chol}(b),(c) gives $T'(Tx)=x$ and $T(T'x)=x$ for all $x$. So $T$ and $T'$ are mutually
inverse bijections; they are linear by Definition~\ref{ar:def:lin2}. By Lemma~\ref{ar:lem:lin2}(b), applied to
$T=\operatorname{lin2}(s,\beta/s,0,t)$, the orientation factor is $st-(\beta/s)\cdot0=st>0$. Finally, by
Lemma~\ref{ar:lem:chol}(a),
$\operatorname{nsq}(Tz)=s^2z_0^2+2\beta z_0z_1+(\beta^2/s^2+t^2)z_1^2$. Here $s^2=\alpha$ and
\[
\frac{\beta^2}{s^2}+t^2=\frac{\beta^2}{\alpha}+\frac{\alpha\gamma-\beta^2}{\alpha}=\gamma .
\]
So $\operatorname{nsq}(Tz)=\alpha z_0^2+2\beta z_0z_1+\gamma z_1^2=\operatorname{avgQ}(L,m,z)$ by
Lemma~\ref{ar:lem:avgq-eq}.

(c) Let $x,y\in\R^2$ and put $w=T'x-T'y$. Since $g(T'x)-g(T'y)=(LT'x+b)-(LT'y+b)=Lw$ and $T$ is linear,
\begin{align*}
\operatorname{eDist2}(h(x),h(y))
&=\operatorname{nsq}\bigl(T(g(T'x))-T(g(T'y))\bigr) &&\text{Lemma~\ref{ar:lem:nsq}(b)}\\
&=\operatorname{nsq}\bigl(T(Lw)\bigr)=\operatorname{avgQ}(L,m,Lw) &&\text{linearity of $T$; part (b)}\\
&=\operatorname{avgQ}(L,m,w)=\operatorname{nsq}(Tw) &&\text{Lemma~\ref{ar:lem:avgq-inv} ($L^m=1$); part (b)}\\
&=\operatorname{nsq}\bigl(T(T'x)-T(T'y)\bigr)=\operatorname{nsq}(x-y) &&\text{linearity of $T$; $TT'=\mathrm{id}$}\\
&=\operatorname{eDist2}(x,y) &&\text{Lemma~\ref{ar:lem:nsq}(b).}
\end{align*}
So $h$ is a Euclidean isometry. If $X\in T(S)$, write $X=Tp$ with $p\in S$. Then $h(X)=T(g(T'(Tp)))=T(g(p))\in T(S)$,
because $g(p)\in S$.

(d) Induction on $k$. For $k=0$: $h^0=\mathrm{id}=T\circ T'=T\circ g^0\circ T'$. If $h^k=T\circ g^k\circ T'$, then
$h^{k+1}=h\circ h^k=T\circ g\circ(T'\circ T)\circ g^k\circ T'=T\circ g^{k+1}\circ T'$, because $T'\circ T=\mathrm{id}$.

If $g^k=\mathrm{id}$, then $h^k=T\circ T'=\mathrm{id}$. If $h^k=\mathrm{id}$, then for every $p$,
$T(g^k(p))=T(g^k(T'(Tp)))=h^k(Tp)=Tp$, so $g^k(p)=p$ because $T$ is injective. The case $k=1$ gives $h\ne\mathrm{id}$
if and only if $g\ne\mathrm{id}$. Since $h^k=\mathrm{id}\Leftrightarrow g^k=\mathrm{id}$ for every $k$, the least
$k\ge1$ with this property is the same for $g$ and $h$.

\emph{Relation to Lean.} Parts (b) and (c) are steps of the Lean proof of
\lean{hasConvexNGon_of_affine_symmetry}, where $T'$ is called \lean{Ti}. These steps are: the positivity of $\alpha$,
of $\alpha\gamma-\beta^2$, of $s$ and of $t$; the inverse identities \lean{hTiT} and \lean{hTTi}, and the injectivity of
$T$ and $T'$ derived from them (\lean{hTinj}, \lean{hTiinj}); the identity \lean{hnT}; the orientation factor
$st-(\beta/s)\cdot0$ of $T$, which Lean uses only inside the step \lean{hgp'} (general position of $T(S)$, through
\lean{generalPosition_image_of_orient} and \lean{orient_lin2}); and the steps \lean{hiso} and \lean{hinv'}. Of
part (d), only the implication $g\ne\mathrm{id}\Rightarrow h\ne\mathrm{id}$ is in Lean (step \lean{hne'}).
\end{proof}

\begin{theorem}[affine-to-Euclidean transfer; \leanref{ErdosSzekeres.SymReduction.hasConvexNGon_of_affine_symmetry}{SYMLEAN/AffineReduction.lean:1330}]\label{ar:thm:transfer}
Let $N,n\in\N$. Assume
\begin{itemize}
\item[\textup{(hE)}] for every finite set $S'\subseteq\R^2$ and every map $h:\R^2\to\R^2$: if $S'$ is in general
position, $N\le\card{S'}$, $h$ is a Euclidean isometry \textup(\lean{IsEucIsometry}~$h$: $\operatorname{eDist2}(h(p),
h(q))=\operatorname{eDist2}(p,q)$ for all $p,q\in\R^2$\textup), $h\ne\mathrm{id}$, and $h(p)\in S'$ for all $p\in S'$,
then \lean{HasConvexNGon}~$n$~$S'$
\end{itemize}
\textup(that is, assume $\mathcal E(N,n)$ of Definition~\textup{\ref{ar:def:E}}\textup), and assume $3\le N$. Let $S$
be a finite set, $g:\R^2\to\R^2$ a map, $L$ a matrix and $b\in\R^2$, such that $g(p)=Lp+b$ for every $p\in\R^2$, $g$ is
injective, $S$ is in general position, $N\le\card S$, $g\ne\mathrm{id}$, and $g(p)\in S$ for every $p\in S$. Then
\lean{HasConvexNGon}~$n$~$S$: some $n$ points of $S$ are in convex position.
\end{theorem}

\begin{proof}
Since $\card S\ge N\ge3$, Proposition~\ref{ar:prop:averaging} applies. Fix $m\ge1$ with $L^m=1$ (part (a)), and let
$T$, $T'$ and $h=T\circ g\circ T'$ be as in parts (b) and (c). We check the hypotheses of (hE) for $S'=T(S)$ and $h$.
\begin{itemize}
\item $T(S)$ is in general position, by Lemma~\ref{ar:lem:gp-orient} with $f=T$ and $D=st\ne0$
(Proposition~\ref{ar:prop:averaging}(b)).
\item $N\le\card S=\card{T(S)}$, because $T$ is injective ($T'\circ T=\mathrm{id}$).
\item $h$ is a Euclidean isometry and maps $T(S)$ into $T(S)$ (Proposition~\ref{ar:prop:averaging}(c)).
\item $h\ne\mathrm{id}$, by Proposition~\ref{ar:prop:averaging}(d), since $g\ne\mathrm{id}$.
\end{itemize}
So (hE) gives $U\subseteq T(S)$ with $\card U=n$ in convex position. Put $V=T'(U)$.
\begin{itemize}
\item $V\subseteq S$: every $y\in V$ is $y=T'(x)$ with $x\in U\subseteq T(S)$, so $x=Tp$ with $p\in S$, and
$y=T'(Tp)=p\in S$.
\item $\card V=\card U=n$, because $T'$ is injective ($T\circ T'=\mathrm{id}$).
\item $V$ is in convex position, by Lemma~\ref{ar:lem:cp-affine} applied to $T'$. This is a linear map, hence an
affine map (Lean \lean{Ti.toAffineMap}), and it is injective.
\end{itemize}
Hence \lean{HasConvexNGon}~$n$~$S$.
\end{proof}

\begin{remark}[the hypothesis (hE), and what the statement forgets]\label{ar:rem:hE}
The following three points concern the hypotheses of Theorem~\ref{ar:thm:transfer} and what its statement does not
record.
\begin{enumerate}[label=(\arabic*)]
\item The hypothesis (hE) is a parameter of Theorem~\ref{ar:thm:transfer}: the Euclidean statement $\mathcal E(N,n)$
for the \emph{same} $N$ and $n$. The theorem is the implication from $\mathcal E(N,n)$ to the corresponding statement
for injective affine maps. It involves no unproved claim of the project.
\item (hE) is used once, for the set $T(S)$ and the isometry $h=T\circ g\circ T'$. Since (hE) quantifies over
\emph{all} Euclidean isometries other than the identity, the statement does not record which isometry $h$ is. The
written Proposition~\ref{ar:prop:averaging}(d) adds that $h$ has the same order as $g$. This is not part of any Lean
statement. An order-preserving form of the transfer, whose Euclidean hypothesis concerns only isometries of the order
of $g$, is not formalized. The remark on affine symmetries in \dirn{WRITEUP/P1/p1.tex} (Remark \texttt{rem:affine})
needs such a form, and says that it is not formalized.
\item The hypothesis $3\le N$ is used only to find three distinct points of $S$. Injectivity of $g$ is a separate
hypothesis. It follows from $\det L\ne0$ (Lemma~\ref{ar:lem:inj}), and that is how
Corollary~\ref{ar:cor:transfer-forms}(c) supplies it.
\end{enumerate}
\end{remark}

\begin{corollary}[the transfer for other presentations of the map; \leanref{ErdosSzekeres.SymReduction.hasConvexNGon_of_affineMap_symmetry}{SYMLEAN/AffineReduction.lean:1423}, \leanref{ErdosSzekeres.SymReduction.hasConvexNGon_of_affineEquiv_symmetry}{SYMLEAN/AffineReduction.lean:1433}, \leanref{ErdosSzekeres.SymReduction.hasConvexNGon_of_affine_coords_symmetry}{SYMLEAN/AffineReduction.lean:1444}]\label{ar:cor:transfer-forms}
Let $N,n\in\N$. Assume \textup{(hE)}, exactly as in Theorem~\textup{\ref{ar:thm:transfer}} \textup(the statement
$\mathcal E(N,n)$\textup), and assume $3\le N$. Let $S$ be a finite set in general position with $N\le\card S$.
\begin{enumerate}[label=\textup{(\alph*)}]
\item \textup{(}\lean{hasConvexNGon_of_affineMap_symmetry}\textup{)} If $f$ is an affine map of $\R^2$ \textup(Mathlib
\lean{AffineMap}\textup) that is injective, $f\ne\mathrm{id}$ as maps, and $f(p)\in S$ for all $p\in S$, then
\lean{HasConvexNGon}~$n$~$S$.
\item \textup{(}\lean{hasConvexNGon_of_affineEquiv_symmetry}\textup{)} If $f$ is an affine equivalence of $\R^2$
\textup(Mathlib \lean{AffineEquiv}\textup) with $f\ne\mathrm{refl}$ and $f(p)\in S$ for all $p\in S$, then
\lean{HasConvexNGon}~$n$~$S$.
\item \textup{(}\lean{hasConvexNGon_of_affine_coords_symmetry}\textup{)} If $g:\R^2\to\R^2$ is a map and
$a,b,c,d,e,f\in\R$ \textup(here $f$ is a real number\textup) satisfy $ad-bc\ne0$,
$g(p)=(ap_0+bp_1+e,\ cp_0+dp_1+f)$ for every $p\in\R^2$, $g\ne\mathrm{id}$, and $g(p)\in S$ for all $p\in S$, then
\lean{HasConvexNGon}~$n$~$S$.
\end{enumerate}
\end{corollary}

\begin{proof}
(a) By Lemma~\ref{ar:lem:matrix-form}, $f(p)=M_fp+f(0)$ for every $p$. Apply Theorem~\ref{ar:thm:transfer} with $g=f$,
$L=M_f$ and $b=f(0)$; the remaining hypotheses are those of (a).

(b) The underlying affine map of $f$ has the same map, which is injective because $f$ is a bijection. It differs from
$\mathrm{id}$: if the map $f$ were $\mathrm{id}$, then $f$ and $\mathrm{refl}$ would have the same map, so
$f=\mathrm{refl}$ by \lean{AffineEquiv.ext}, a contradiction. Now apply (a).

(c) Let $L=\begin{pmatrix}a&b\\c&d\end{pmatrix}$ and $b'=(e,f)$. By Lemma~\ref{ar:lem:mulvec},
$Lp+b'=(ap_0+bp_1+e,\ cp_0+dp_1+f)=g(p)$ for every $p$. Moreover $\det L=ad-bc\ne0$ (Mathlib
\lean{Matrix.det_fin_two_of}), so $g$ is injective by Lemma~\ref{ar:lem:inj}. Apply Theorem~\ref{ar:thm:transfer}.
\end{proof}

\subsection{Affine symmetries and convex heptagons}

\begin{theorem}[convex heptagons from affine symmetries, conditional on the five case classes; \leanref{ErdosSzekeres.SymReduction.heptagon_of_nontrivial_affineEquiv}{SYMLEAN/AffineReduction.lean:1501}, \leanref{ErdosSzekeres.SymReduction.heptagon_of_nontrivial_affineEquiv_image}{SYMLEAN/AffineReduction.lean:1510}, \leanref{ErdosSzekeres.SymReduction.heptagon_of_nontrivial_affineMap}{SYMLEAN/AffineReduction.lean:1489}, \leanref{ErdosSzekeres.SymReduction.heptagon_of_nontrivial_affine}{SYMLEAN/AffineReduction.lean:1477}, \leanref{ErdosSzekeres.SymReduction.heptagon_of_nontrivial_affine_coords}{SYMLEAN/AffineReduction.lean:1521}]\label{ar:thm:main}
Assume, as hypotheses, the five propositions \lean{HalfTurn34}, \lean{Threefold33}, \lean{Fivefold35},
\lean{MirrorOneAxis33} and \lean{MirrorNoAxis34} of Definition~\textup{\ref{sr:def:classes}}. Let $S\subseteq\R^2$ be a
finite set in general position with $33\le\card S$.
\begin{enumerate}[label=\textup{(\alph*)}]
\item \textup{(}\lean{heptagon_of_nontrivial_affineEquiv}\textup{)} If $g$ is an affine equivalence of $\R^2$
\textup(Mathlib \lean{AffineEquiv}\textup) with $g\ne\mathrm{refl}$ and $g(p)\in S$ for all $p\in S$, then
\lean{HasConvexNGon}~$7$~$S$.
\item \textup{(}\lean{heptagon_of_nontrivial_affineEquiv_image}\textup{)} If $g$ is an affine equivalence of $\R^2$
with $g\ne\mathrm{refl}$ and $g(S)\subseteq S$ \textup(an inclusion of sets: the image of $S$ under $g$ is contained in
$S$\textup), then \lean{HasConvexNGon}~$7$~$S$.
\item \textup{(}\lean{heptagon_of_nontrivial_affineMap}\textup{)} If $f$ is an affine map of $\R^2$ \textup(Mathlib
\lean{AffineMap}\textup) that is injective, $f\ne\mathrm{id}$ as maps, and $f(p)\in S$ for all $p\in S$, then
\lean{HasConvexNGon}~$7$~$S$.
\item \textup{(}\lean{heptagon_of_nontrivial_affine}\textup{)} If $g:\R^2\to\R^2$ is a map, $L$ a matrix and
$b\in\R^2$, with $g(p)=Lp+b$ for every $p$, $g$ injective, $g\ne\mathrm{id}$, and $g(p)\in S$ for all $p\in S$, then
\lean{HasConvexNGon}~$7$~$S$.
\item \textup{(}\lean{heptagon_of_nontrivial_affine_coords}\textup{)} If $g:\R^2\to\R^2$ is a map and
$a,b,c,d,e,f\in\R$ \textup(in this part $b$ and $f$ are real numbers\textup) satisfy $ad-bc\ne0$,
$g(p)=(ap_0+bp_1+e,\ cp_0+dp_1+f)$ for every $p\in\R^2$, $g\ne\mathrm{id}$, and $g(p)\in S$ for all $p\in S$, then
\lean{HasConvexNGon}~$7$~$S$.
\end{enumerate}
In each case, \lean{HasConvexNGon}~$7$~$S$ means that some $7$ points of $S$ are in convex position.
\end{theorem}

\begin{proof}
By Theorem~\ref{sr:thm:main} (\lean{heptagon_of_nontrivial_isometry}, applied to the five hypotheses), the following
holds. \emph{For every finite set $S'$ and every map $h:\R^2\to\R^2$: if $S'$ is in general position, $33\le\card{S'}$,
$h$ is a Euclidean isometry, $h\ne\mathrm{id}$ and $h$ maps $S'$ into $S'$, then \lean{HasConvexNGon}~$7$~$S'$.} This is
exactly $\mathcal E(33,7)$, the hypothesis (hE) with $N=33$ and $n=7$. Moreover $3\le33$.

(a) Apply Corollary~\ref{ar:cor:transfer-forms}(b).

(b) For $p\in S$ we have $g(p)\in g(S)\subseteq S$, so $g$ maps $S$ into $S$. Now apply (a).

(c) Apply Corollary~\ref{ar:cor:transfer-forms}(a).

(d) Apply Theorem~\ref{ar:thm:transfer}.

(e) Apply Corollary~\ref{ar:cor:transfer-forms}(c).

The Lean proofs are the same applications, each a two-line term. For (a), (c), (d) and (e) they pass
\lean{heptagon_of_nontrivial_isometry}, applied to the five hypotheses \lean{h2}, \lean{h3}, \lean{h5}, \lean{hm1} and
\lean{hm0}, as \lean{hE}, and prove $3\le33$ by \lean{norm_num}. For (b) the Lean proof passes no \lean{hE}: it applies
\lean{heptagon_of_nontrivial_affineEquiv}, that is (a), after turning the inclusion into maps-into.
\end{proof}

\begin{remark}[the five hypotheses]\label{ar:rem:classes}
Theorem~\ref{ar:thm:main} is a conditional statement. Its five hypotheses enter only through
Theorem~\ref{sr:thm:main}. In Lean they are definitions of type \lean{Prop} that occur as explicit premises
(\lean{h2}, \lean{h3}, \lean{h5}, \lean{hm1}, \lean{hm0}), not axioms. In short, each class says that a finite set in
general position contains $7$ points in convex position if it has one of the following: exactly $34$ points and a
half-turn symmetry whose centre is not in the set; exactly $33$ points and a symmetry under the rotation by $2\pi/3$
about a centre not in the set; exactly $35$ points and a symmetry under the rotation by $2\pi/5$ about a centre not in
the set; exactly $33$ points, a reflection symmetry in a line, and exactly one point of the set on that line; exactly
$34$ points, a reflection symmetry, and no point of the set on the line. See Definition~\ref{sr:def:classes} for the
exact statements. In the project they are supplied as follows (\dirn{SYMLEAN/README.md}, \S\S5, 6 and~9;
\dirn{WRITEUP/SYMMETRIC.md}, \S\S1, 3 and~4).
\begin{itemize}
\item \lean{MirrorOneAxis33} and \lean{MirrorNoAxis34} are Lean-checked theorems: \lean{mirrorOneAxis33_proved}
(line 4274) and \lean{mirrorNoAxis34_proved} (line 4282) of \dirn{SYMLEAN/mirror/ReductionThreeClasses.lean}
(Theorem~\ref{mb:thm:classes}). These two theorems are about the copies of the class definitions in that file. It does
not import \dirn{SYMLEAN/AffineReduction.lean}; lines 38--1559 of \dirn{SYMLEAN/mirror/ReductionThreeClasses.lean} are
byte-identical to lines 12--1533 of \dirn{SYMLEAN/AffineReduction.lean}. So the two theorems discharge the premises
\lean{hm1} and \lean{hm0} of Theorem~\ref{ar:thm:main} only through the identity of the two texts. Their axioms log
\dirn{SYMLEAN/mirror/axioms.out} lists only \texttt{propext}, \texttt{Classical.choice} and \texttt{Quot.sound}. They
are obtained without SAT input from the repository's Lean formalization of Baek's theorem
(Theorem~\ref{mb:thm:baek}) and the mirror reduction. The same file restates part (a) with
only the three rotation classes as hypotheses (\lean{heptagon_of_nontrivial_affineEquiv_3}, line 4303;
Theorem~\ref{mb:thm:main3}). Both mirror classes also have SAT certificates (record \dirn{SYMZ}).
\item \lean{HalfTurn34}, \lean{Threefold33} and \lean{Fivefold35} are not Lean-checked. They are computer-assisted
claims of the records \dirn{C2FINAL} (with \dirn{C2C34} and \dirn{C2P}--\dirn{C2P5}), \dirn{C3FINAL} and \dirn{SYM5}.
They rest on four kinds of evidence.
\begin{itemize}
\item \emph{SAT certificates.} Every certificate that these claims use, including those of the pieces of their case
splits, carries a recorded acceptance by the proof checker \texttt{cake\_lpr}. Since a revision of 6 November 2026 this is the
only checker in the project's trust base; the second checker, \texttt{lrat-check}, is a secondary check only
(\dirn{WRITEUP/SYMMETRIC.md}, \S3; \dirn{SYMLEAN/README.md}, \S9, which predates that revision, still names both). For
the parallelogram pieces of \lean{HalfTurn34} and the $116{,}883$ pieces of \lean{Threefold33}, the proof files were
deleted after checking, so for most of these pieces only the recorded verdicts remain. Some of these pieces were
re-solved (\dirn{C2FINAL}, \S\S4, 5 and~9; \dirn{C2RECHECK}; \dirn{C3SPOT}; \dirn{C3PIECES}).
\item \emph{Written proofs, reviewed or audited within the project} (listed in \dirn{WRITEUP/SYMMETRIC.md}, \S3).
Among them are the soundness proofs of the encoders; the reductions A and B of the parallelogram cases of
\lean{HalfTurn34}; the normal form of \lean{Threefold33} (\dirn{C3-S/soundness.md}) and Lemma~5 and Lemma~H6M of its
case split; and, for \lean{Threefold33} and \lean{Fivefold35}, a perturbation argument for centres that lie on a line
through two points of different orbits.
\item \emph{Exact computations} that compose the pieces (\dirn{C2FINAL}, \dirn{C3FINAL}). Two bookkeeping steps are
computer-assisted because their finite checks are kernel evaluations (\texttt{decide +kernel}): the cover of the $910$
threefold roots, which is Lean-checked (Theorem~\ref{cb:thm:c3}), and the closure of the central parallelogram orbits,
which is Lean-checked according to the record (\dirn{C2FINAL/README.md}, \S6; Theorem~\ref{cb:thm:c2}). Lean does not
see the records: that every leaf named by the keys of the certificates of Theorem~\ref{cb:thm:c3}
and every subcase listed in the data of Theorem~\ref{cb:thm:c2} has a record with the checker verdicts and a CNF hash equal to that of the rebuilt formula is an exact computation outside
Lean.
\item \emph{The match between the classes and the formulas.} No formal link joins a Lean class to its CNF formulas.
That each class states what its certificates prove was checked by reading (\dirn{SYMLEAN/README.md}, \S6).
\end{itemize}
The records list the remaining caveats and open checks (\dirn{WRITEUP/SYMMETRIC.md}, \S4).
\end{itemize}
The unconditional statement, Theorem~\ref{ar:thm:main} without its five hypotheses, is claimed in the project as
computer-assisted (Theorem~A of \dirn{WRITEUP/SYMMETRIC.md}, \S0, in its form for affine bijections), with the evidence
listed above; it is not Lean-checked as a whole. This document establishes it only conditionally: as
Theorem~\ref{ar:thm:main} and, for part~(a) with three hypotheses, as Theorem~\ref{mb:thm:main3}.
\end{remark}

\begin{remark}[what the statements of this section do not say]\label{ar:rem:scope}
The statements of this section are limited in four ways.
\begin{enumerate}[label=(\arabic*)]
\item The symmetry is a map of the whole plane: an affine equivalence, an injective affine map, or an injective map
$p\mapsto Lp+b$. A permutation of $S$ that is not the restriction of such a map is not treated.
\item One symmetry is treated at a time. Nothing is stated about the full symmetry group of $S$.
\item The variant without \lean{HalfTurn34} (Theorem~\ref{sr:thm:variant}) has no affine version in the Lean files
(\dirn{SYMLEAN/README.md}, \S9).
\item The order of the symmetry is not tracked (Remark~\ref{ar:rem:hE}(2)).
\end{enumerate}
\end{remark}

\paragraph{Coverage.}
This section covers every declaration of \dirn{SYMLEAN/AffineReduction.lean}, lines 1059--1533:
\begin{itemize}
\item \lean{nsq} (Definition~\ref{ar:def:nsq}); \lean{mulVec2_fst} and \lean{mulVec2_snd}
(Lemma~\ref{ar:lem:mulvec}); \lean{nsq_nonneg} and \lean{eDist2_eq_nsq} (Lemma~\ref{ar:lem:nsq});
\item \lean{gA}, \lean{gB} and \lean{gC} (Definition~\ref{ar:def:gram}); \lean{nsq_mulVec}
(Lemma~\ref{ar:lem:nsq-mulvec});
\item \lean{avgQ}, \lean{qA}, \lean{qB} and \lean{qC} (Definition~\ref{ar:def:avgq}); \lean{avgQ_eq}
(Lemma~\ref{ar:lem:avgq-eq}); \lean{avgQ_ge}, \lean{qA_pos} and \lean{qdisc_pos} (Lemma~\ref{ar:lem:avgq-pos});
\lean{avgQ_mulVec} (Lemma~\ref{ar:lem:avgq-inv});
\item \lean{lin2} (Definition~\ref{ar:def:lin2}); \lean{lin2_fst}, \lean{lin2_snd} and \lean{orient_lin2}
(Lemma~\ref{ar:lem:lin2}); \lean{nsq_lin2_chol}, \lean{lin2_chol_inv_left} and \lean{lin2_chol_inv_right}
(Lemma~\ref{ar:lem:chol});
\item \lean{iterate_sub_of_affine} (Lemma~\ref{ar:lem:iter}); \lean{exists_periodic_of_injective}
(Lemma~\ref{ar:lem:periodic}); \lean{matrix_eq_one_of_fix} (Lemma~\ref{ar:lem:fix}); \lean{orient_affine}
(Lemma~\ref{ar:lem:orient-affine}); \lean{injective_of_affine} (Lemma~\ref{ar:lem:inj}); \lean{affineMap_eq_mulVec}
(Lemma~\ref{ar:lem:matrix-form});
\item \lean{generalPosition_image_of_orient} (Lemma~\ref{ar:lem:gp-orient}); \lean{generalPosition_image_affine} and
\lean{generalPosition_image_affineEquiv} (Lemma~\ref{ar:lem:gp-affine}); \lean{convexPosition_image_affine}
(Lemma~\ref{ar:lem:cp-affine}); \lean{hasConvexNGon_image_affine} (Lemma~\ref{ar:lem:ngon-affine});
\item \lean{hasConvexNGon_of_affine_symmetry} (Theorem~\ref{ar:thm:transfer}, conditional on (hE));
\lean{hasConvexNGon_of_affineMap_symmetry}, \lean{hasConvexNGon_of_affineEquiv_symmetry} and
\lean{hasConvexNGon_of_affine_coords_symmetry} (Corollary~\ref{ar:cor:transfer-forms}, conditional on (hE));
\item \lean{heptagon_of_nontrivial_affineEquiv}, \lean{heptagon_of_nontrivial_affineEquiv_image},
\lean{heptagon_of_nontrivial_affineMap}, \lean{heptagon_of_nontrivial_affine} and
\lean{heptagon_of_nontrivial_affine_coords} (Theorem~\ref{ar:thm:main}, conditional on the case classes
\lean{HalfTurn34}, \lean{Threefold33}, \lean{Fivefold35}, \lean{MirrorOneAxis33} and \lean{MirrorNoAxis34}).
\end{itemize}
Definition~\ref{ar:def:E} and Proposition~\ref{ar:prop:averaging} are written steps, not Lean declarations.

Deliberately not covered here:
\begin{itemize}
\item Lines 1--1058 of \dirn{SYMLEAN/AffineReduction.lean}. They are a verbatim duplicate of lines 1--1058 of
\dirn{SYMLEAN/SymmetricReduction.lean} and are treated in the section of Theorem~\ref{sr:thm:main}. Of them, \lean{iterate_mem}
(line 69) is re-derived inside the proof of Lemma~\ref{ar:lem:periodic}, and \lean{point_ext} (line 502: two points
with equal coordinates are equal) is used tacitly.
\item \dirn{SYMLEAN/AffineReductionAxioms.lean}. It is the same file followed by a blank line, nine
\texttt{\#print axioms} lines and one \texttt{\#check} line, and it is used here only as evidence for the axioms.
\item The verbatim copy of lines 12--1533 inside \dirn{SYMLEAN/mirror/ReductionThreeClasses.lean} (the declarations of
this section sit there at their line number plus 26). The own content of that file is treated
in the section of Theorem~\ref{mb:thm:main3}.
\end{itemize}
\endgroup
